% ============================================================================
% PLANTILLA 07 — CONTROL DE LECTURA
% ----------------------------------------------------------------------------
% Para: controles de lectura, lectura crítica, comentario de textos,
% interpretación de teorías y ensayos breves con rúbrica.
% Compilar: latexmk -pdf plantilla-07-control-de-lectura.tex
% Solucionario: \documentclass[soluciones]{academic-exam}
% ============================================================================
\documentclass{academic-exam}

\universidad{Universidad Nacional de San Cristóbal de Huamanga}
\facultad{Facultad de Ciencias Económicas, Administrativas y Contables}
\escuela{Escuela Profesional de Economía}
\curso{Economía Política}
\codigocurso{EC-121}
\docente{Mg. Nombre Apellido}
\periodo{Semestre 2026-I}
\tipoevaluacion{Control de Lectura 03}
\fechaevaluacion{21 de julio de 2026}
\duracion{80 minutos}
\modalidad{Presencial · individual}

\begin{document}
\membrete

\begin{datos}[Lectura evaluada]
\textsc{Polanyi, Karl} (1944). \emph{La gran transformación: los orígenes
políticos y económicos de nuestro tiempo}. Capítulos 4--6:
\enquote{Las sociedades y los sistemas económicos}, \enquote{La evolución
del modelo de mercado} y \enquote{El mercado autorregulado y las
mercancías ficticias}.
\end{datos}

\begin{instrucciones}
\begin{itemize}
  \item Responda con sus propias palabras; se penaliza la transcripción
        literal del texto sin elaboración.
  \item Cite pasajes de la lectura solo como apoyo de su argumento,
        indicando el capítulo.
  \item Extensión sugerida por respuesta: entre media y una carilla.
\end{itemize}
\end{instrucciones}

% ============================================================ PREGUNTAS ====
\begin{pregunta}[4][Comprensión]
Explique qué entiende Polanyi por \emph{mercancías ficticias} y por qué
sostiene que la tierra, el trabajo y el dinero pertenecen a esa
categoría.
\lineasrespuesta{9}
\begin{solucion}
Una respuesta completa identifica: las mercancías ficticias son elementos
que no fueron \emph{producidos para la venta} — condición que define a la
mercancía genuina —, pero que el sistema de mercado trata como si lo
fueran. El trabajo es actividad humana inseparable de la vida; la tierra
es la naturaleza no producida por el hombre; el dinero es un símbolo del
poder de compra creado institucionalmente. Someterlos por completo al
mecanismo de precios amenaza con destruir la sociedad y el entorno
natural, lo que engendra el \emph{doble movimiento} de protección social.
\end{solucion}
\end{pregunta}

\begin{pregunta}[4][Interpretación contextual]
El autor escribe:
\begin{displayquote}
  Permitir que el mecanismo de mercado sea el único director del destino
  de los seres humanos y de su medio natural [...] tendría como resultado
  la destrucción de la sociedad.
\end{displayquote}
Interprete este pasaje a la luz del contexto histórico en el que fue
escrito (Inglaterra industrial y crisis de entreguerras).
\lineasrespuesta{9}
\begin{solucion}
Se espera que el estudiante vincule el pasaje con: (i) la experiencia de
la ley de pobres y Speenhamland en la Inglaterra del siglo XIX;
(ii) el colapso del patrón oro y la rigidez deflacionaria de los años
veinte; (iii) el ascenso de los fascismos como reacción autoritaria al
desarraigo del mercado autorregulado. La tesis central: el mercado
autorregulado es una utopía institucional cuyo intento de realización
plena produce contramovimientos defensivos de la sociedad.
\end{solucion}
\end{pregunta}

\begin{pregunta}[5][Análisis crítico]
Polanyi afirma que la economía estuvo históricamente \emph{incrustada}
(\emph{embedded}) en relaciones sociales, y que el siglo XIX invirtió esa
relación. Contraste esta tesis con la visión de la economía neoclásica
sobre el mercado. ¿Qué gana y qué pierde el análisis económico con cada
enfoque?
\lineasrespuesta{11}
\begin{solucion}
Respuesta de nivel excelente: la teoría neoclásica modela el mercado como
mecanismo de asignación universal y ahistórico, con agentes maximizadores
e instituciones exógenas; gana precisión formal, capacidad predictiva en
mercados competitivos y estática comparativa. El enfoque polanyiano gana
profundidad histórica e institucional — reciprocidad, redistribución e
intercambio como formas de integración — y explica por qué los mercados
requieren armazón estatal, pero pierde capacidad de formalización y de
predicción puntual. Se valora el uso de ejemplos (mercado laboral,
regulación financiera).
\end{solucion}
\end{pregunta}

\begin{pregunta}[7][Comentario crítico]
Redacte un comentario crítico (máximo dos carillas) sobre la vigencia de
la tesis del \emph{doble movimiento} para analizar la economía peruana
contemporánea. Elija un caso concreto (informalidad laboral, conflictos
socioambientales mineros, regulación financiera u otro) y organice su
texto con introducción, desarrollo y conclusión.

\begin{criterios}[Rúbrica del comentario crítico]
\small
\renewcommand{\arraystretch}{1.35}
\begin{tabularx}{\linewidth}{@{}p{2.9cm} X X c@{}}
  \textbf{Criterio} & \textbf{Logrado} & \textbf{En proceso} & \textbf{Pje.}\\
  \arrayrulecolor{grislinea}\hline\addlinespace[3pt]
  Tesis y estructura &
    Posición clara; texto organizado en intro\-ducción, desarrollo y cierre. &
    Posición implícita o estructura incompleta. & 2\\
  Uso de la lectura &
    Conceptos de Polanyi empleados con precisión y citas pertinentes. &
    Menciones genéricas sin precisión conceptual. & 2\\
  Aplicación al caso &
    El caso peruano se analiza con evidencia concreta. &
    Caso descrito sin conexión analítica. & 2\\
  Redacción &
    Prosa clara, cohesionada y sin errores. &
    Errores frecuentes que dificultan la lectura. & 1\\
\end{tabularx}
\end{criterios}
\lineasrespuesta{10}
\end{pregunta}

\findelexamen
\end{document}
